\section{Síntesis y ampliación}

\subsection{Trama de tecnologías centrales}

Esta lección parte de la pregunta básica de la generación condicional e introduce sistemáticamente el stack tecnológico clave para llevar los modelos de difusión a aplicaciones prácticas:

\begin{table}[H]
\centering
\begin{tabular}{lp{4.5cm}p{4.5cm}}
\toprule
\textbf{Tecnología} & \textbf{Problema resuelto} & \textbf{Idea central} \\
\midrule
Classifier Guidance & Generación condicional & score = score incondicional + gradiente del clasificador \\
CFG & Eliminación del clasificador externo & dropout condicional + extrapolación en inferencia \\
DDIM Inversion & Edición de imágenes & inversión determinista + regeneración cambiando condiciones \\
Latent Diffusion & Reducción de costos computacionales & compresión previa al espacio latente antes de la difusión \\
ControlNet & Control condicional con pocos datos & modelo base congelado + copia entrenable del codificador \\
LoRA & Personalización eficiente & actualización de pesos de rango bajo \\
\bottomrule
\end{tabular}
\caption{Resumen de las tecnologías centrales de esta lección}
\end{table}

\subsection{Perspectiva unificada: Manipulación de puntajes (Score Manipulation)}

Una perspectiva que atraviesa toda esta lección (y todo el campo de los modelos de difusión) es la \textbf{manipulación de puntajes} —dirigir las trayectorias de generación modificando la función de puntaje (o la predicción de ruido equivalente), guiándolas hacia una dirección objetivo. Ya sea Classifier Guidance, CFG o más aplicaciones que se discutirán en lecciones futuras, su esencia es:

$$
\tilde{s}(x_t, t) = s_{\text{base}}(x_t, t) + \text{término de guía}
$$

Esta es la ventaja central que distingue a los modelos de difusión de GAN y VAE: cada paso del proceso generativo puede ser ajustado flexiblemente mediante señales externas.

\subsection{Recomendaciones de elección de ingeniería: cuándo usar qué tecnología}

Si se aplican los contenidos de esta lección al diseño de sistemas reales, la siguiente tabla permite juzgar rápidamente la selección tecnológica:

\begin{table}[H]
\centering
\begin{tabular}{p{0.2\textwidth} p{0.24\textwidth} p{0.26\textwidth}}
\toprule
Objetivo de tarea & Tecnología prioritaria & Criterios de decisión clave \\
\midrule
Solo se requiere generación condicional con texto & CFG / Latent Diffusion & ¿Se acepta cierto costo computacional adicional a cambio de una consistencia condicional más fuerte? \\
Se requiere control estructural (bordes, postura, profundidad) & ControlNet & ¿Ya existe una señal condicional confiable y si se busca adaptación rápida con pocos datos \\
Se requiere estilo personalizado o personalización del sujeto & LoRA & ¿Hay solo少量 de datos? ¿Se exige baja memoria gráfica y despliegue combinable \\
Se requiere conservar la estructura de la imagen original y editarla & DDIM Inversion + regeneración cambiando condiciones & ¿Se exige modificar el contenido de la imagen original de forma controlada en lugar de muestrear desde cero? \\
\bottomrule
\end{tabular}
\caption{Tabla de referencia rápida para selección de tecnologías en Lecture 6}
\end{table}

\subsection{Significado para las lecciones subsiguientes}

El valor de Lecture 6 no radica solo en presentar varios métodos específicos, sino en establecer un marco unificado: los modelos de difusión preentrenados son un prior general, mientras que CFG, ControlNet, LoRA e Inversion redirigen ese prior hacia tareas específicas a diferentes niveles. Ya sea guía durante inferencia, difusión para video o generación 3D, todo esto esencialmente se puede ver como una extensión adicional de este marco.

\subsection{Lecturas adicionales}

\begin{itemize}
    \item \textbf{Classifier Guidance}: Dhariwal \& Nichol, ``Diffusion Models Beat GANs on Image Synthesis,'' NeurIPS 2021
    \item \textbf{Classifier-Free Guidance}: Ho \& Salimans, ``Classifier-Free Diffusion Guidance,'' NeurIPS Workshop 2022
    \item \textbf{Latent Diffusion / Stable Diffusion}: Rombach et al., ``High-Resolution Image Synthesis with Latent Diffusion Models,'' CVPR 2022
    \item \textbf{ControlNet}: Zhang et al., ``Adding Conditional Control to Text-to-Image Diffusion Models,'' ICCV 2023
    \item \textbf{LoRA}: Hu et al., ``LoRA: Low-Rank Adaptation of Large Language Models,'' ICLR 2022
    \item \textbf{DDIM}: Song et al., ``Denoising Diffusion Implicit Models,'' ICLR 2021
    \item \textbf{HuggingFace Diffusers}: \href{https://huggingface.co/docs/diffusers}{huggingface.co/docs/diffusers} — biblioteca estándar de modelos de difusión preentrenados
\end{itemize}

%% --- Fin del contenido --- %%

\end{document}
